Voy a resolver cada integral doble en el orden más conveniente.  

---



---

### **Inciso (b)**  

La región \( D \) está delimitada por \( y = x \) y \( y = x^3 \).  

#### **Límites de integración:**  
- **Tipo II:** \( 0 \leq y \leq 1 \), \( y^{1/3} \leq x \leq y \).  

Resolvemos en **Tipo II**:  

\[
I = \int_0^1 \int_{y^{1/3}}^y (x^2 + 2y) \, dx \, dy.
\]

La integral interna:

\[
\int (x^2 + 2y) dx = \frac{x^3}{3} + 2yx.
\]

Evaluamos en \( x = y \) y \( x = y^{1/3} \):

\[
\left( \frac{y^3}{3} + 2y^2 \right) - \left( \frac{y}{3} + 2y^{4/3} \right).
\]

\[
I = \int_0^1 \left( \frac{y^3}{3} + 2y^2 - \frac{y}{3} - 2y^{4/3} \right) dy.
\]

Integrando:

\[
\frac{y^4}{12} + \frac{2y^3}{3} - \frac{y^2}{6} - \frac{6y^{7/3}}{7} \Big|_0^1.
\]

\[
I = \left( \frac{1}{12} + \frac{2}{3} - \frac{1}{6} - \frac{6}{7} \right).
\]

Llevando a común denominador:

\[
I = \frac{7}{84} + \frac{56}{84} - \frac{14}{84} - \frac{72}{84} = \frac{-23}{84}.
\]

---

### **Inciso (c)**  

La región \( D \) es el triángulo con vértices \( (0,1) \), \( (1,2) \), \( (4,1) \).  

#### **Ecuaciones de las rectas:**  
- \( y = x + 1 \) (de \( (0,1) \) a \( (1,2) \)).  
- \( y = -\frac{1}{3} x + \frac{4}{3} \) (de \( (1,2) \) a \( (4,1) \)).  

#### **Límites de integración en Tipo I:**  
- Para \( 0 \leq x \leq 1 \), \( 1 \leq y \leq x + 1 \).  
- Para \( 1 \leq x \leq 4 \), \( -\frac{1}{3} x + \frac{4}{3} \leq y \leq 1 \).  

Resolvemos:

\[
I = \int_0^1 \int_1^{x+1} \ln(y) \, dy \, dx + \int_1^4 \int_{-\frac{1}{3}x+\frac{4}{3}}^1 \ln(y) \, dy \, dx.
\]

Integrando:

\[
\int \ln(y) dy = y \ln(y) - y.
\]

Evaluando en los límites y sumando ambas integrales, obtenemos:

\[
I = \frac{5}{6} \ln 2 + \frac{5}{6} - \frac{1}{2} \ln 3 - \frac{1}{2}.
\]

---

### **Inciso (d)**  

La región \( D \) está acotada por \( x = 0 \) y \( x = \sqrt{1 - y^2} \), es decir, es la parte derecha del círculo \( x^2 + y^2 = 1 \).  

#### **Límites de integración:**  
- **Tipo II:** \( -1 \leq y \leq 1 \), \( 0 \leq x \leq \sqrt{1 - y^2} \).  

Resolvemos en **Tipo II**:  

\[
I = \int_{-1}^{1} \int_0^{\sqrt{1 - y^2}} x y^2 \, dx \, dy.
\]

La integral interna:

\[
\int x y^2 dx = \frac{x^2}{2} y^2.
\]

Evaluando:

\[
I = \int_{-1}^{1} \frac{(1 - y^2) y^2}{2} dy.
\]

\[
= \frac{1}{2} \int_{-1}^{1} (y^2 - y^4) dy.
\]

Calculamos:

\[
\frac{1}{2} \left( \frac{y^3}{3} - \frac{y^5}{5} \right) \Big|_{-1}^{1}.
\]

\[
= \frac{1}{2} \left( \frac{1}{3} - \frac{1}{5} - (-\frac{1}{3} + \frac{1}{5}) \right).
\]

\[
= \frac{1}{2} \left( \frac{2}{15} - (-\frac{2}{15}) \right).
\]

\[
= \frac{1}{2} \left( \frac{4}{15} \right) = \frac{2}{15}.
\]

---

### **Respuestas Finales**  
a) \( \frac{1 - \cos 1}{2} \).  
b) \( -\frac{23}{84} \).  
c) \( \frac{5}{6} \ln 2 + \frac{5}{6} - \frac{1}{2} \ln 3 - \frac{1}{2} \).  
d) \( \frac{2}{15} \).